\documentclass[11pt]{article}
\usepackage[margin=1in]{geometry}
\usepackage[T1]{fontenc}
\usepackage{amsmath,amssymb}
\usepackage{booktabs}
\usepackage{graphicx}
\usepackage{url}
\usepackage{caption}
\captionsetup{font=small,labelfont=bf}
\renewcommand{\thetable}{S\arabic{table}}
\renewcommand{\thefigure}{S\arabic{figure}}
\renewcommand{\thesection}{S\arabic{section}}
\setlength{\parskip}{0.4em}
\title{\vspace{-2em}Supporting Information\\[0.4em]
\large A testing-effort artefact in the most-sensitive aquatic endpoint:
structure-based PFAS screening recovers predictive validity but still cannot
support clearance}
\author{Jeongyun Kim, Jin Yoo, Ku Kang}
\date{}
\begin{document}
\maketitle

\section{Extraction from the ECOTOX Knowledgebase}

All toxicity data derive from the U.S. EPA ECOTOX Knowledgebase ASCII release
\texttt{ecotox\_ascii\_09\_15\_2026}, obtained from
\url{https://gaftp.epa.gov/ecotox/}. The release contains 725{,}636 test records
and 1{,}250{,}612 result records. We used it directly rather than through a
curated derivative, so that the number of records per substance --- the quantity
at issue in this paper --- reflects the primary compilation.

Selection proceeded in the order shown in Table~S1. Tests were restricted to
aquatic organisms in freshwater media; results were restricted to acute LC50,
EC50 and IC50 and chronic NOEC, LOEC, LOEL, NOEL and MATC endpoints reported in
mass-concentration units; and values carrying an inequality qualifier were
separated out. Records in molar units (\texttt{uM}, \texttt{nM},
\texttt{mM}, \texttt{nmol/L}, \texttt{umol/L}) and in dose or body-weight units
(\texttt{ppm diet}, \texttt{mg/kg bdwt}) were not converted and were excluded;
molar units alone account for 3{,}967 of the rejected PFAS records.

\begin{table}[h]
\centering
\caption{Record selection for the PFAS set.}
\small
\begin{tabular}{lr}
\toprule
Step & Count \\
\midrule
Tests on PFASSTRUCTV6 CAS numbers & 11{,}719 \\
\quad after restricting to aquatic organisms in freshwater & 3{,}816 \\
Results joined to those tests & 29{,}375 \\
\quad with an acute or chronic endpoint of interest & 24{,}379 \\
\quad with convertible mass-concentration units and a numeric value & 16{,}357 \\
\quad carrying an inequality qualifier (excluded) & 366 \\
\textbf{Usable uncensored records} & \textbf{15{,}991} \\
Unique substances represented & 103 \\
\bottomrule
\end{tabular}
\end{table}

The endpoint composition of the PFAS set is strongly asymmetric
(Table~S2) and this matters for the argument of the paper. Chronic sublethal
endpoints outnumber acute endpoints by roughly fifteen to one, and they also lie
lower on the concentration scale. A protective minimum taken over all strata is
therefore drawn almost always from the chronic sublethal pool, and the size of
that pool is what varies most between well- and poorly-studied substances.

\begin{table}[h]
\centering
\caption{Endpoint composition of the 15{,}991 usable PFAS records.}
\small
\begin{tabular}{lrlr}
\toprule
Endpoint & Records & Endpoint & Records \\
\midrule
NOEC & 8{,}584 & LC50 & 437 \\
LOEC & 6{,}360 & NOEL & 43 \\
EC50 & 523    & IC50 & 18 \\
MATC & 17     & LOEL & 9 \\
\bottomrule
\end{tabular}
\end{table}

The same filters applied without the PFAS restriction, and limited to acute
endpoints, gave 106{,}095 freshwater acute records for 4{,}757 chemicals, which
is the set used for the cross-chemical replication in Section~S2 and as the
training pool in Section~S5.

\section{Endpoint constructions and the effort diagnostic}

Four summaries of the same records were compared, all expressed as
$y=\log_{10}(T/\mathrm{mg\,L^{-1}})$: the minimum over all strata (the
protective convention), the median over the same records, the median over acute
records only, and a depth-equalised minimum described in Section~S3. Testing
effort was taken as $\log_{10}$ of the number of usable records per substance.

\begin{table}[h]
\centering
\caption{Association between each endpoint and testing effort. No model is
involved; this is a property of the response variable.}
\small
\begin{tabular}{llrr}
\toprule
Set & Endpoint & $n$ substances & Spearman $\rho$ vs $\log_{10}$ records \\
\midrule
PFAS & minimum (all strata) & 103 & $-0.734$ \\
PFAS & median (all strata)  & 103 & $-0.442$ \\
PFAS & acute median         & 70  & $-0.036$ \\
\midrule
All chemicals & acute minimum & 4{,}757 & $-0.462$ \\
All chemicals & acute median  & 4{,}757 & $-0.160$ \\
\bottomrule
\end{tabular}
\end{table}

The last two rows are the cross-chemical replication. Restricting to acute data
does not by itself remove the confound: within acute records alone, the minimum
still carries roughly three times the association of the median. The confound
attaches to the act of taking a minimum, not to the inclusion of chronic data,
although the two compound in the PFAS set because the chronic pool is both
larger and lower.

\begin{table}[h]
\centering
\caption{Protective minimum by testing depth, PFAS set.}
\small
\begin{tabular}{lrrrr}
\toprule
Records per substance & $n$ & Median minimum & Median spread & Median of medians \\
\midrule
1--3     & 26 & $+1.40$ & $0.00$ & $+1.46$ \\
4--10    & 15 & $-0.23$ & $1.90$ & $+0.72$ \\
11--50   & 26 & $-1.60$ & $2.00$ & $+0.24$ \\
51--300  & 26 & $-2.73$ & $5.09$ & $-0.82$ \\
$>$300   & 10 & $-4.96$ & $7.22$ & $-1.16$ \\
\bottomrule
\end{tabular}
\end{table}

Spread is the within-substance range in $\log_{10}$ units and rises with depth
as expected for a sample range. The final column shows that the median is far
less sensitive to depth than the minimum, which is the basis for preferring it.

\section{Equal-depth subsampling}

To separate the extreme-order-statistic artefact from genuine reverse
causation, every substance with at least $k$ records was subsampled to exactly
$k$ records without replacement, the minimum recomputed, and the correlation
with the original record count recalculated. The procedure was repeated 200
times per $k$.

\begin{table}[h]
\centering
\caption{Correlation between the protective minimum and testing effort after
equalising sampling depth. Ranges are 2.5th--97.5th percentiles over 200 draws.}
\small
\begin{tabular}{lrrl}
\toprule
$k$ & Eligible substances & Mean $\rho$ & 95\% range \\
\midrule
all records & 103 & $-0.734$ & --- \\
3  & 78 & $-0.311$ & $[-0.412,-0.207]$ \\
5  & 70 & $-0.279$ & $[-0.387,-0.163]$ \\
10 & 64 & $-0.293$ & $[-0.373,-0.204]$ \\
20 & 47 & $-0.524$ & $[-0.600,-0.435]$ \\
\bottomrule
\end{tabular}
\end{table}

At $k=3$, $5$ and $10$ the correlation stabilises near $-0.29$, slightly under
half the value obtained using all records. The $k=20$ row is higher, but its 47
eligible substances are by construction the well-studied ones, so the comparison
is made over a restricted and shifted range; we report it for completeness and
do not treat it as evidence against the pattern.

\section{Fixed-model comparison across endpoint constructions}

In this experiment the broad non-PFAS model of the published analysis was held
fixed and only the external target was changed, so that differences cannot be
attributed to training, tuning or representation.

\begin{table}[h]
\centering
\caption{External performance of one fixed model against four constructions of
the same underlying records. Confidence intervals are substance-level bootstrap
(4{,}000 resamples) on Spearman $\rho$.}
\small
\begin{tabular}{lrrrlr}
\toprule
Endpoint & $n$ & $R^2$ & $\rho$ & 95\% CI on $\rho$ & $\rho$ vs effort \\
\midrule
minimum (published)      & 103 & $-0.745$ & $-0.015$ & $[-0.21,+0.18]$ & $-0.734$ \\
median                   & 103 & $-0.147$ & $+0.076$ & $[-0.13,+0.27]$ & $-0.442$ \\
depth-equalised minimum  & 70  & $-2.053$ & $+0.208$ & $[-0.03,+0.41]$ & $-0.332$ \\
acute median             & 70  & $+0.216$ & $+0.424$ & $[+0.21,+0.61]$ & $-0.036$ \\
\bottomrule
\end{tabular}
\end{table}

The ordering of external performance follows the ordering of freedom from the
effort confound. The depth-equalised minimum is an informative intermediate: its
rank correlation improves over the raw minimum, consistent with removing part of
the confound, but its $R^2$ is very poor because averaging subsampled minima
produces a target with compressed variance that the model is not calibrated to.
We report it as a diagnostic rather than as a candidate endpoint.

\section{Consistent rebuild}

Training and external sets were rebuilt on the acute median. Model settings were
those of the published analysis and were not re-tuned: 2{,}048-bit ECFP4
fingerprints (Morgan radius 2), random forest with 300 trees,
\texttt{max\_features} $=\sqrt{p}$, seed 20260620. Broad cross-validation used
five repeats of shuffled five-fold splitting with 160 trees per fold. Conformal
validation used 40 repeats of the normalised split-conformal procedure given in
the main text, with marginal calibration on a random 80/20 split of the broad
set and PFAS-conditional calibration on disjoint halves of the PFAS set.

\begin{table}[h]
\centering
\caption{Published protective endpoint versus the consistent acute-median
rebuild. Model, representation and conformal procedure are identical.}
\small
\begin{tabular}{lrr}
\toprule
Quantity & Protective minimum & Acute median \\
\midrule
Broad training structures & 3{,}788 & 3{,}718 \\
External PFAS structures & 119 & 58 \\
Broad $5\times5$ CV $R^2$ & $0.387\pm0.024$ & $0.443\pm0.029$ \\
External $R^2$ & $-0.580$ & $+0.367$ \\
External Spearman $\rho$ & $-0.071$ & $+0.460$ \\
\quad 95\% CI & --- & $[+0.22,+0.65]$ \\
Marginal coverage at nominal 0.90 & 0.659 & 0.858 \\
Marginal width & 4.84 & 3.70 \\
PFAS-conditional coverage & 0.922 & 0.897 \\
PFAS-conditional width & 10.53 & 4.46 \\
Inventory structures in domain & 17{,}128 & 16{,}070 \\
\bottomrule
\end{tabular}
\end{table}

The external set is smaller on the acute median for two reasons: restricting to
acute endpoints removes substances with only chronic data, and charge-neutralised
desalting collapses salt forms of the same parent. Both reductions are
deliberate, and the smaller external set is the principal limitation of the
rebuilt analysis.

\section{Clearance across hazard thresholds}

Because the two endpoints differ in protectiveness, a single threshold does not
provide a comparable decision rule. On the 70 substances with both summaries the
acute median lies $+2.69$ $\log_{10}$ units above the protective minimum, and
94.3\% of measured PFAS exceed 0.1 mg/L on the acute median against 42.9\% on
the minimum. Clearance was therefore evaluated across a continuum.

\begin{table}[h]
\centering
\caption{Clearance of in-domain inventory structures ($n=16{,}070$) under the
rebuilt model, by hazard threshold. The required half-width is the median over
in-domain structures of the distance from the point prediction to the threshold;
the achieved median half-width is 2.44.}
\small
\begin{tabular}{lrrr}
\toprule
Threshold & Cleared & \% & Median required half-width \\
\midrule
0.01 mg/L & 14{,}991 & 93.3 & $+3.22$ \\
0.1 mg/L (used with the protective endpoint) & 5{,}441 & 33.9 & $+2.22$ \\
1 mg/L (GHS acute aquatic Category 1) & 306 & 1.9 & $+1.22$ \\
10 mg/L & 0 & 0.0 & $+0.22$ \\
49 mg/L (matched to the protective endpoint) & 0 & 0.0 & $-0.47$ \\
100 mg/L & 0 & 0.0 & $-0.78$ \\
\bottomrule
\end{tabular}
\end{table}

Two regimes are visible. Above 10 mg/L the required half-width approaches and
then falls below zero, meaning the median in-domain structure is predicted more
hazardous than the threshold and no narrowing of the interval would clear it.
Below that, the barrier is interval width alone: at the GHS Category 1 boundary
the median structure would clear with a 90\% interval of 2.43 $\log_{10}$ units
against the 4.87 achieved.

\section{Structure standardisation and resolution}

ECOTOX identifies substances by CAS number and carries no structures. CAS
numbers were converted to the dashed form and resolved to SMILES through the
PubChem PUG REST interface at the published rate limit of five requests per
second; 4{,}311 of 4{,}757 chemicals resolved (90.6\%), and the remainder were
dropped. Resolution failures are not random with respect to chemical class ---
mixtures, poorly defined substances and some organometallics resolve less often
--- and we were not able to characterise the resulting bias.

Two standardisation choices differ from common practice and are stated
explicitly because they change the data.

\emph{Desalting keeps the fluorine-bearing fragment, not the largest one.} The
usual rule retains the largest fragment by heavy-atom count. For
tetrabutylammonium perfluoro-1-butanesulfonate the tetrabutylammonium cation
(17 heavy atoms) exceeds the perfluorobutanesulfonate anion (13 heavy atoms), so
the largest-fragment rule discards the analyte and models the counterion. We
therefore retain the fluorinated fragment where one exists, falling back to the
largest fragment otherwise.

\emph{Fragments are charge-neutralised before canonicalisation.} Without this
step a free acid and its salts produce different canonical SMILES and are
counted as distinct structures. In the measured PFAS set, desalting alone
collapses 138 rows to 131 parents, whereas desalting with neutralisation
collapses them to 119. The difference is not cosmetic: perfluorooctanesulfonate
appears in seven rows whose assigned protective values span 7.93 $\log_{10}$
units, and without neutralisation six of those rows survive as separate
substances and can fall on both sides of a calibration split.

\section{OECD subclass assignment and statistical power}

Substances were assigned to structure categories following the OECD (2021)
definition and the 2018 OECD/UNEP structure series. The definitional gate
requires at least one fully fluorinated, saturated methyl or methylene carbon
carrying no hydrogen, chlorine, bromine or iodine. Carbons that are fully
halogenated but not fully fluorinated (CHF$_2$, CClF$_2$), methine carbons
($>$CF--) alone, and unsaturated carbons ($=$CF$_2$) do not satisfy it.

Applying the gate to the measured set excluded seven substances, listed here
because several would be retained by a looser fluorine-count rule:
4-bromo-3-chloro-3,4,4-trifluorobut-1-ene and
1,1,2-trichloro-1,2,2-trifluoroethane (halogens other than fluorine on the
candidate carbon); perfluorobiphenyl, 2,3,4,5,6-pentafluoroaniline,
decafluorobenzophenone and pentafluorobenzaldehyde (fluorine attached only to
aromatic rings); and difluoroacetic acid, whose CHF$_2$ carbon bears a hydrogen.

Categories were assigned by parsing the molecule as C$_n$F$_{2n+1}$--X--R with
head-group identity taking precedence over incidental fluorination. Assignment
uses connectivity from fluorinated carbons rather than requiring a literal
CF$_2$ bonded to the head group, so that perfluoroalkyl ether acids
(for example HFPO-DA, whose $\alpha$-carbon is a methine) and perfluorinated
cycloalkane sulfonic acids are classified by head group rather than falling
through to a residual category.

\begin{table}[h]
\centering
\caption{External rank correlation by OECD structure class on the protective
endpoint, where the measured set is largest. Intervals are substance-level
bootstrap.}
\small
\begin{tabular}{lrrl}
\toprule
Class & $n$ & Spearman $\rho$ & 95\% CI \\
\midrule
All OECD-2021 PFAS & 112 & $+0.002$ & $[-0.18,+0.20]$ \\
\midrule
200a perfluoroalkane sulfonic & 11 & $-0.611$ & $[-0.92,+0.03]$ \\
100 perfluoroalkyl carboxylic & 15 & $-0.206$ & $[-0.74,+0.47]$ \\
700 side-chain fluorinated aromatic & 12 & $-0.007$ & $[-0.52,+0.61]$ \\
400 fluorotelomer & 23 & $+0.274$ & $[-0.20,+0.63]$ \\
200b perfluoroalkane sulfonamide & 11 & $+0.323$ & $[-0.41,+0.87]$ \\
600 other perfluoroalkyl & 17 & $+0.350$ & $[-0.25,+0.71]$ \\
500 perfluoroalkyl ether acid & 13 & $+0.695$ & $[+0.14,+0.93]$ \\
\bottomrule
\end{tabular}
\end{table}

Six of seven intervals include zero. The seventh is one of seven comparisons
and rests on 13 substances; at that sample size 80\% power detects only
$|\rho|\ge0.71$, so a point estimate of $+0.695$ is approximately the smallest
value the design could have returned as significant, and is the configuration
most exposed to selection among the seven.

On the acute-median set the position is worse. Of 68 OECD-classified
substances, three classes reach ten members and the largest reaches fourteen.

\begin{table}[h]
\centering
\caption{Detectable effect size by class on the acute-median set, at 80\% power
and $\alpha=0.05$.}
\small
\begin{tabular}{lrr}
\toprule
Class & $n$ & Detectable $|\rho|$ \\
\midrule
100 perfluoroalkyl carboxylic & 14 & 0.69 \\
200a perfluoroalkane sulfonic & 14 & 0.69 \\
400 fluorotelomer & 11 & 0.76 \\
700 side-chain fluorinated aromatic & 9 & 0.82 \\
600 other perfluoroalkyl & 7 & 0.89 \\
200b perfluoroalkane sulfonamide & 6 & 0.92 \\
\bottomrule
\end{tabular}
\end{table}

\begin{table}[h]
\centering
\caption{Measured substances required per class to resolve differential transfer,
at 80\% power.}
\small
\begin{tabular}{lrr}
\toprule
Target effect size & $\alpha=0.05$ & $\alpha=0.05/7$ (Bonferroni) \\
\midrule
$\rho=0.3$ & 85 & 133 \\
$\rho=0.4$ & 47 & 73 \\
$\rho=0.5$ & 29 & 44 \\
$\rho=0.6$ & 19 & 29 \\
\bottomrule
\end{tabular}
\end{table}

Resolving class-level differences at $\rho=0.3$ across seven classes would
require roughly 600 measured substances uncorrected, or 930 with correction,
against the 103 the public freshwater record currently supports. We report this
as the measurement requirement implied by the question rather than as an
argument that the question is unanswerable.

\section{Files supporting the numerical audit}

\begin{itemize}\setlength{\itemsep}{0.1em}
\item \texttt{extract.py} --- PFAS record extraction and per-substance summaries.
\item \texttt{effort\_test.py} --- effort correlations and equal-depth subsampling (Tables~S3--S5).
\item \texttt{build\_broad.py} --- acute-median candidate set across all 4{,}757 chemicals.
\item \texttt{resolve\_smiles.py} --- CAS to SMILES resolution via PubChem, resume-safe.
\item \texttt{endpoint\_fix.py} --- fixed-model comparison across endpoint constructions (Table~S6).
\item \texttt{final\_pipeline.py} --- consistent rebuild: cross-validation, external validation, conformal intervals, clearance (Tables~S7--S8).
\item \texttt{oecd\_classify.py} --- OECD (2021) definitional gate and category assignment (Tables~S9--S11).
\item \texttt{make\_fig\_effort.py}, \texttt{make\_fig\_final.py}, \texttt{make\_forest.py} --- Figures 1--3.
\item \texttt{ecotox\_pfas\_measured.csv}, \texttt{ecotox\_acute\_smiles.csv} --- per-substance summaries with record counts.
\item \texttt{endpoint\_fix\_results.csv}, \texttt{effort\_subsample.csv}, \texttt{final\_metrics.json} --- values in Tables~S5--S8.
\end{itemize}

\end{document}
